\section{Methodology}

Our analysis combines three feature families---lexical, stylometric, and semantic---spanning linguistic abstraction from exact n-gram overlap (most sensitive to any change) to contextual embeddings (most robust to rephrasing), so relative decay rates reveal which levels of representation are most fragile.

\subsection{Feature Extraction}

\textbf{Stylometric features.} We extract Part-of-Speech tags using spaCy \cite{honnibal2020spacy} (\texttt{en\_core\_web\_sm}), counting the 17 Universal POS tags per text and normalising to proportions. This yields a length-invariant 17-dim vector capturing the syntactic ``skeleton'' independent of word choice.

We also extract named entity sets via spaCy NER (PERSON, ORG, GPE, DATE, etc.), lowercased to avoid case mismatches. Named entities are the semantic ``cargo''; unlike distributional features, entity preservation requires \emph{exact} lexical reproduction, making it a much sharper probe of transformation.

We also extract dependency tree depth via spaCy's parser: the max root-to-leaf depth per sentence, averaged across the document. This is a proxy for syntactic complexity---deeper trees indicate more subordination.

\textbf{Lexical metrics.} We compute BLEU \cite{papineni2002bleu} (sentence-level smoothed n-gram precision) and ROUGE-L \cite{lin2004rouge} (longest common subsequence F1) between $T_0$ and each paraphrased version. These measure surface word overlap and sequential preservation respectively. We emphasize that these are purely lexical metrics: a BLEU of 0.0 does not mean the text has lost its meaning, only that the specific words have been replaced.

\textbf{Semantic metrics.} We compute BERTScore \cite{zhang2020bertscore} with DeBERTa-xlarge-MNLI, comparing texts in contextual embedding space. High BERTScore with low BLEU indicates meaning is preserved despite lexical replacement---the signature of successful paraphrasing. We prefer BERTScore to SBERT cosine because token-level F1 from an NLI-tuned encoder correlates more tightly with human adequacy on paraphrase perturbations. Type--Token Ratio is omitted as a standalone feature: it is collinear with document length, and word count plus the 17-dim POS vector already carry the same lexical-diversity signal at higher resolution.

\subsection{Similarity Analysis (RQ1)}

To measure style vs.\ content decay, we employ two complementary signals at the structural level:

\textbf{POS cosine similarity:} For each article, we compute the cosine similarity between the $T_0$ POS proportion vector and the corresponding $T_n$ vector ($n \in \{1,2,3\}$). Cosine similarity is appropriate because POS vectors are compositional proportions: we care about the \emph{direction} (relative tag distribution) rather than the magnitude. A score near 1.0 indicates the grammatical distribution is unchanged; lower scores indicate syntactic restructuring. We report mean $\pm$ standard deviation across all articles, broken down by dataset and paraphraser.

\textbf{NER entity metrics:} For each article, we compute three set-based metrics between $T_0$ and $T_n$ entity sets:
\begin{itemize}
    \item \textbf{Jaccard similarity} $J = |A \cap B| / |A \cup B|$: overall entity overlap.
    \item \textbf{Recall} $R = |A \cap B| / |A|$: fraction of original entities retained (low $R$ = entity dropping).
    \item \textbf{Precision} $P = |A \cap B| / |B|$: fraction of $T_n$ entities that existed in $T_0$ (low $P$ = novel entity introduction).
\end{itemize}
Separating Recall and Precision is critical because dropping and novel introduction are distinct phenomena with different implications: dropping indicates lossy compression, while novel introduction indicates the paraphraser is substituting entities. Articles where $T_0$ contains zero entities are excluded from means to avoid undefined metrics.

\subsection{Authorship Attribution (RQ2)}

To locate the ``point of no return'' at which stylometric identity becomes unrecoverable, we frame authorship attribution as a binary classification task (Human vs.\ LLM source) over the 23-dimensional feature vector extracted in \S4.1 (17 POS proportions, dependency depth, word count, sentence count, punctuation density, type-token ratio proxy, and entity count). We train three classifiers with complementary inductive biases---Logistic Regression (linear, interpretable), SVM with RBF kernel (non-linear margin), and Random Forest (axis-aligned ensemble)---each with balanced class weighting and standardised features. Critically, all three are trained \emph{only} on $T_0$ features and then evaluated, unchanged, on the same articles at $T_1$, $T_2$, and $T_3$. This protocol mirrors a realistic deployment: a forensic classifier is calibrated on known-clean text and must hold up under progressive adversarial rewriting. Macro F1 is reported because the Human vs.\ LLM classes are moderately imbalanced across the seven source authors, and the random baseline of 0.5 provides a clean ``signal lost'' threshold.

\subsection{Paraphraser Identification (RQ3)}

For RQ3 we move from attribution to \emph{tool identification}: given a paraphrased text, can we recover which paraphraser produced it? We construct 24-dimensional delta feature vectors by subtracting $T_0$ features from $T_1$ features (17 POS-proportion deltas, dependency-depth delta, word-count ratio, sentence-count ratio, plus NER Jaccard, Recall, and Precision between $T_0$ and $T_1$ entity sets). Using deltas rather than raw $T_1$ features cancels the per-document authorial baseline and isolates the transformation signature of the paraphraser itself. We train a 6-class classifier (SVM with RBF kernel, matched against Logistic Regression and Random Forest for comparison). PaLM2 is excluded because only 37 document keys have paraphrases from all seven paraphrasers; excluding PaLM2 yields 1,546 common keys. Splits follow the document-group protocol of \S3.3 so that a document never appears in both train and test sets, and macro F1 is the headline metric. A confusion matrix and Random Forest feature importance are reported for interpretability.

\subsection{Controlled Experiments}

We design two controlled comparisons to isolate specific variables, holding all other factors constant:

\textbf{Dipper intensity:} Dipper \cite{krishna2023dipper} offers Low and High intensity settings with identical architecture. Comparing POS cosine vs.\ NER Recall across these settings on the same articles isolates the effect of paraphrasing aggressiveness on style vs.\ content (RQ1).

\textbf{Pegasus coverage:} Pegasus (Slight) \cite{zhang2020pegasus} paraphrases 25\% of sentences; Pegasus (Full) paraphrases all of them. Same model, different coverage, isolating whether broader coverage primarily drops entities or substitutes them.

\textbf{Source author comparison:} We split by whether $T_0$ was human-written or produced by one of six LLMs and compare POS drift, testing whether the original author's syntactic patterns affect how much a paraphraser restructures.

\subsection{Data Alignment}

For the stylometric analysis we build a forensic pivot with one row per (document key, source author, dataset), and columns for $T_0$ and the seven $T_1$ paraphraser variants. Version strings are parsed with a regex tokenizer that handles parenthesised tokens (e.g., \texttt{dipper(low)\_dipper(low)}). After dropping rows with missing columns, the aligned corpus contains \textbf{18,058 article-source pairs}; POS vectors, NER sets, and dependency depths are cached to disk.

For the multi-tier analysis ($T_0 \rightarrow T_1 \rightarrow T_2 \rightarrow T_3$) we build separate per-paraphraser pivots, yielding 508 (PaLM2) to 3,018 (Pegasus Slight) matched chains. Metrics are always computed relative to $T_0$ to capture cumulative decay.

\subsection{Statistical Validation}

All reported means are accompanied by 95\% bootstrap confidence intervals (1,000 resamples, percentile method). RQ2 classifiers are additionally reported with 5-fold stratified cross-validation on $T_0$ (mean $\pm$ std macro F1) before being retrained on the full $T_0$ set for the $T_1$--$T_3$ decay curves. RQ3 classifiers are compared against three trivial baselines (random uniform, most-frequent, stratified) to contextualise the learned performance. Key comparisons between paraphrasers and between tiers are validated with two-sided paired permutation tests (10,000 permutations).
